% GENERATED FILE. Edit canonical lessons, then run npm run generate:books.
% canonical-source-hash: fnv1a64:068d9561fc26fa14
% canonical-lessons: FR-C362-citoyen, FR-C362-citoyennete, FR-C362-nationalite, FR-C362-liberte, FR-C362-egalite

\chapter{Citoyens}
\label{ch:fr-citoyens}
\hlchaptermodality{362}

\begin{chapteropening}
\textbf{By the end of this chapter:} \emph{I can talk about citizens, citizenship, nationality, freedom and equality.}

Complete the last lesson of chapter 362: I can talk about citizens, citizenship, nationality, freedom and equality.
\end{chapteropening}

\section[le citoyen]{le citoyen — the citizen}

\label{lesson:FR-C362-citoyen}



\begin{quote}
Before the new one: say the French for the efficiency, then the French for the productivity.
\end{quote}

\subsection*{You'll want to know: le citoyen}
\textbf{le citoyen} — \textquotedblleft{}the citizen\textquotedblright{}.

\begin{quote}
Chaque citoyen a le droit de voter.
\end{quote}

\begin{quote}
\emph{Every citizen has the right to vote.}
\end{quote}

Feminine \emph{citoyenne}.

\begin{cousinweb}
\emph{citoyen} is inherited from Old French \emph{citoien} (EtymDB), built on \emph{cité} \textquotedblleft{}city\textquotedblright{}, from Latin \emph{cīvitās} \textquotedblleft{}citizenship, a state\textquotedblright{}, from \emph{cīvis} \textquotedblleft{}citizen\textquotedblright{} (Lewis \& Short).
\end{cousinweb}

\begin{grammarlens}[title={What you've built}]
One more word for telling, explaining and describing at B1.
\end{grammarlens}

\subsection*{Your turn}
\begin{itemize}
\raggedright
  \item \emph{Say it:} \emph{le citoyen}
  \item \emph{Say it:} \emph{le citoyen}, once more
  \item \emph{Say it:} a sentence of your own with \emph{le citoyen}
  \item \emph{Recall:} say the French for the efficiency, then the French for the productivity, then say \emph{le citoyen} again
\end{itemize}

\begin{tcolorbox}[breakable,colback=teal!4,colframe=teal!35!black,title={Before you move on}]
What does le citoyen mean? (\textquotedblleft{}The citizen\textquotedblright{}.) Say it once more.
\end{tcolorbox}
\section[la citoyenneté]{la citoyenneté — the citizenship}

\label{lesson:FR-C362-citoyennete}



\begin{quote}
Before the new one: say the French for the productivity, then the French for the citizen.
\end{quote}

\subsection*{You'll want to know: la citoyenneté}
\textbf{la citoyenneté} — \textquotedblleft{}the citizenship\textquotedblright{}.

\begin{quote}
Il a demandé la citoyenneté de son nouveau pays.
\end{quote}

\begin{quote}
\emph{He applied for citizenship of his new country.}
\end{quote}

From \emph{citoyen} + \emph{-eté}.

\begin{cousinweb}
\emph{citoyenneté} is built from \emph{citoyen} \textquotedblleft{}citizen\textquotedblright{} and \emph{-eté}. \emph{citoyen} goes back to Latin \emph{cīvitās} \textquotedblleft{}citizenship\textquotedblright{}, from \emph{cīvis} \textquotedblleft{}citizen\textquotedblright{} (Lewis \& Short).
\end{cousinweb}

\begin{grammarlens}[title={What you've built}]
One more word for telling, explaining and describing at B1.
\end{grammarlens}

\subsection*{Your turn}
\begin{itemize}
\raggedright
  \item \emph{Say it:} \emph{la citoyenneté}
  \item \emph{Say it:} \emph{la citoyenneté}, once more
  \item \emph{Say it:} a sentence of your own with \emph{la citoyenneté}
  \item \emph{Recall:} say the French for the productivity, then the French for the citizen, then say \emph{la citoyenneté} again
\end{itemize}

\begin{tcolorbox}[breakable,colback=teal!4,colframe=teal!35!black,title={Before you move on}]
What does la citoyenneté mean? (\textquotedblleft{}The citizenship\textquotedblright{}.) Say it once more.
\end{tcolorbox}
\section[la nationalité]{la nationalité — the nationality}

\label{lesson:FR-C362-nationalite}



\begin{quote}
Before the new one: say the French for the citizen, then the French for the citizenship.
\end{quote}

\subsection*{You'll want to know: la nationalité}
\textbf{la nationalité} — \textquotedblleft{}the nationality\textquotedblright{}.

\begin{quote}
Elle a la double nationalité, française et canadienne.
\end{quote}

\begin{quote}
\emph{She has dual nationality, French and Canadian.}
\end{quote}

\emph{De nationalité française} — of French nationality.

\begin{cousinweb}
\emph{nationalité} is built from \emph{nation} and \emph{-ité}. \emph{nation} comes from Latin \emph{nātiō} \textquotedblleft{}birth, a people\textquotedblright{} (Lewis \& Short), from \emph{nāscī} \textquotedblleft{}to be born\textquotedblright{}.
\end{cousinweb}

\begin{grammarlens}[title={What you've built}]
One more word for telling, explaining and describing at B1.
\end{grammarlens}

\subsection*{Your turn}
\begin{itemize}
\raggedright
  \item \emph{Say it:} \emph{la nationalité}
  \item \emph{Say it:} \emph{la nationalité}, once more
  \item \emph{Say it:} a sentence of your own with \emph{la nationalité}
  \item \emph{Recall:} say the French for the citizen, then the French for the citizenship, then say \emph{la nationalité} again
\end{itemize}

\begin{tcolorbox}[breakable,colback=teal!4,colframe=teal!35!black,title={Before you move on}]
What does la nationalité mean? (\textquotedblleft{}The nationality\textquotedblright{}.) Say it once more.
\end{tcolorbox}
\section[la liberté]{la liberté — the freedom, liberty}

\label{lesson:FR-C362-liberte}



\begin{quote}
Before the new one: say the French for the citizenship, then the French for the nationality.
\end{quote}

\subsection*{You'll want to know: la liberté}
\textbf{la liberté} — \textquotedblleft{}the freedom, liberty\textquotedblright{}.

\begin{quote}
Ces gens se battent pour leur liberté.
\end{quote}

\begin{quote}
\emph{These people are fighting for their freedom.}
\end{quote}

\emph{En liberté} — free, at large.

\begin{cousinweb}
\emph{liberté} comes through Old French \emph{liberté} from Latin \emph{lībertās} \textquotedblleft{}freedom\textquotedblright{} (EtymDB), from \emph{līber} \textquotedblleft{}free\textquotedblright{} (Lewis \& Short).
\end{cousinweb}

\begin{grammarlens}[title={What you've built}]
One more word for telling, explaining and describing at B1.
\end{grammarlens}

\subsection*{Your turn}
\begin{itemize}
\raggedright
  \item \emph{Say it:} \emph{la liberté}
  \item \emph{Say it:} \emph{la liberté}, once more
  \item \emph{Say it:} a sentence of your own with \emph{la liberté}
  \item \emph{Recall:} say the French for the citizenship, then the French for the nationality, then say \emph{la liberté} again
\end{itemize}

\begin{tcolorbox}[breakable,colback=teal!4,colframe=teal!35!black,title={Before you move on}]
What does la liberté mean? (\textquotedblleft{}The freedom, liberty\textquotedblright{}.) Say it once more.
\end{tcolorbox}
\section[l'égalité]{l'égalité — the equality}

\label{lesson:FR-C362-egalite}



\begin{quote}
Before the new one: say the French for the nationality, then the French for the freedom.
\end{quote}

\subsection*{You'll want to know: l'égalité}
\textbf{l'égalité} — \textquotedblleft{}the equality\textquotedblright{}.

\begin{quote}
Ils demandent l'égalité entre les hommes et les femmes.
\end{quote}

\begin{quote}
\emph{They are asking for equality between men and women.}
\end{quote}

In sport, \emph{à égalité} — level, a draw. Feminine.

\begin{cousinweb}
\emph{égalité} is built from \emph{égal} \textquotedblleft{}equal\textquotedblright{} and \emph{-ité} (EtymDB). \emph{égal} comes from Latin \emph{aequālis} \textquotedblleft{}equal\textquotedblright{}, from \emph{aequus} \textquotedblleft{}level, even\textquotedblright{} (Lewis \& Short).
\end{cousinweb}

\begin{grammarlens}[title={What you've built}]
One more word for telling, explaining and describing at B1.
\end{grammarlens}

\subsection*{Your turn}
\begin{itemize}
\raggedright
  \item \emph{Say it:} \emph{l'égalité}
  \item \emph{Say it:} \emph{l'égalité}, once more
  \item \emph{Say it:} a sentence of your own with \emph{l'égalité}
  \item \emph{Recall:} say the French for the nationality, then the French for the freedom, then say \emph{l'égalité} again
\end{itemize}

\begin{tcolorbox}[breakable,colback=teal!4,colframe=teal!35!black,title={Before you move on}]
What does l'égalité mean? (\textquotedblleft{}The equality\textquotedblright{}.) Say it once more.
\end{tcolorbox}
